% Options for packages loaded elsewhere
\PassOptionsToPackage{unicode}{hyperref}
\PassOptionsToPackage{hyphens}{url}
\documentclass[
  11pt,
]{article}
\usepackage{xcolor}
\usepackage[margin=1in]{geometry}
\usepackage{amsmath,amssymb}
\setcounter{secnumdepth}{3}
\usepackage{iftex}
\ifPDFTeX
  \usepackage[T1]{fontenc}
  \usepackage[utf8]{inputenc}
  \usepackage{textcomp} % provide euro and other symbols
\else % if luatex or xetex
  \usepackage{unicode-math} % this also loads fontspec
  \defaultfontfeatures{Scale=MatchLowercase}
  \defaultfontfeatures[\rmfamily]{Ligatures=TeX,Scale=1}
\fi
\usepackage{lmodern}
\ifPDFTeX\else
  % xetex/luatex font selection
\fi
% Use upquote if available, for straight quotes in verbatim environments
\IfFileExists{upquote.sty}{\usepackage{upquote}}{}
\IfFileExists{microtype.sty}{% use microtype if available
  \usepackage[]{microtype}
  \UseMicrotypeSet[protrusion]{basicmath} % disable protrusion for tt fonts
}{}
\makeatletter
\@ifundefined{KOMAClassName}{% if non-KOMA class
  \IfFileExists{parskip.sty}{%
    \usepackage{parskip}
  }{% else
    \setlength{\parindent}{0pt}
    \setlength{\parskip}{6pt plus 2pt minus 1pt}}
}{% if KOMA class
  \KOMAoptions{parskip=half}}
\makeatother
\usepackage{longtable,booktabs,array}
\newcounter{none} % for unnumbered tables
\usepackage{calc} % for calculating minipage widths
% Correct order of tables after \paragraph or \subparagraph
\usepackage{etoolbox}
\makeatletter
\patchcmd\longtable{\par}{\if@noskipsec\mbox{}\fi\par}{}{}
\makeatother
% Allow footnotes in longtable head/foot
\IfFileExists{footnotehyper.sty}{\usepackage{footnotehyper}}{\usepackage{footnote}}
\makesavenoteenv{longtable}
\setlength{\emergencystretch}{3em} % prevent overfull lines
\providecommand{\tightlist}{%
  \setlength{\itemsep}{0pt}\setlength{\parskip}{0pt}}
\usepackage{bookmark}
\IfFileExists{xurl.sty}{\usepackage{xurl}}{} % add URL line breaks if available
\urlstyle{same}
\hypersetup{
  hidelinks,
  pdfcreator={LaTeX via pandoc}
}

\newtheorem{theorem}{Theorem}
\newtheorem{lemma}{Lemma}
\newtheorem{proposition}{Proposition}
\newtheorem{corollary}{Corollary}
\newtheorem{definition}{Definition}
\newtheorem{remark}{Remark}

\author{}
\date{}

\title{Conditional Reduction to General Relativity from Zero-Sum Constraints: Lovelock Uniqueness and the Group-Theoretic Origin of Dimension}

\begin{document}

\maketitle

\begin{abstract}

This work establishes a conditional reduction chain from a discrete
ontological substrate to general relativity. The method does not
assemble the Einstein field equations piece by piece; instead, it
compresses the field equations into three structural
properties---locality \((L)\), second-order \((O)\), and conserved
source \((C')\)---and then invokes the Lovelock uniqueness theorem to
force them into Einstein form. The algebraic skeleton: zero-sum
constraints fix codimension 1; the incidence matrix of a connected graph
spans the zero-sum hyperplane and gives integer reachability; discrete
Dirichlet energy is positive semidefinite with kernel the constants;
continuum approximation combined with metric-volume consistency recovers
the spatial metric; reversible/irreversible splitting directly gives
signature \((1,m-1)\). The metric construction: closed-walk counting
gives a parameter-free edge weight \(w_{ij}\propto\phi_i\phi_jA_{ij}\);
finite truncation \(k=n\) restores locality. The dimension account is
the weakest link: the core postulates do not constrain dimension
(proposition: for every \(m\ge2\) there exists a model satisfying the
postulates); \(D=4\) follows from two independent routes---physically,
``propagating gravitons exist'' requires the unique positive solution of
\(\frac{D(D-3)}2=2\), and geometrically, the closure--motion
correspondence principle \(\dim\ker B=\dim\operatorname{im}B\) combined
with edge-transitivity uniquely selects \(K_4\). The premises of both
routes are not consequences of the core postulates; dimension is
therefore a conditional input.

\end{abstract}

\noindent\textbf{Keywords}: Lovelock theorem; general relativity; zero-sum
constraint; discrete geometry; Hodge duality; dimension problem

\begin{center}\rule{0.5\linewidth}{0.5pt}\end{center}

\section{Introduction}\label{introduction}

The spacetime dimension \(D=4\) is an input rather than a conclusion
within the standard physical framework {[}1,2{]}. Our approach starts
from a discrete substrate containing no manifold, coordinates, metric,
probability, or action, such that the shape of the Einstein field
equations becomes a consequence derivable from the Lovelock uniqueness
theorem.

\textbf{Strategy.} We do not assemble GR piece by piece. GR is
compressed into three structural properties: \((L)\) locality, \((O)\)
second-order, \((C')\) conserved source. The substrate need only provide
these three; the shape of the geometric field equations is then locked
by the Lovelock theorem {[}3{]}.

\textbf{Scope.} The Lovelock theorem's literal statement is ``in 4
dimensions.'' This lever therefore cannot be used to select \(D=4\);
dimension must be argued separately (§6).

\begin{center}\rule{0.5\linewidth}{0.5pt}\end{center}

\section{Core Postulates of the
Substrate}\label{core-postulates-of-the-substrate}

The substrate declares only the following; no manifold, coordinates,
dimension, distance, curvature, action, or probability weights appear:

\begin{itemize}
\tightlist
\item
  \textbf{State}: codimension-1 affine constraint on integer lattice
  (zero-sum);
\item
  \textbf{Transport}: connected, reversible compensatory moves;
\item
  \textbf{Energy}: Dirichlet-type quadratic form;
\item
  \textbf{Time}: step index \(\tau\) and terminal;
\item
  \textbf{Record}: multiplicity multiset of closed classes;
\item
  \textbf{Weights}: integer multiplicities.
\end{itemize}

The substrate axiom is \(Z_0=(\{\pm1\},\iota)+(\mathcal W_n,\mu)\),
where \(\iota\) is the direction involution and \(\mu\) the uniform
counting measure {[}4{]}.

\begin{center}\rule{0.5\linewidth}{0.5pt}\end{center}

\section{Algebraic Skeleton}\label{algebraic-skeleton}

\subsection{Zero-Sum Codimension and Transport
Rank}\label{zero-sum-codimension-and-transport-rank}

\textbf{Lemma 3.1 (Codimension 1).} The zero-sum constraint
\(\sum_v x_v=0\) gives only one linear condition:

\[
\dim H_Q=m-1, \qquad \{e_j-e_1\}_{j=2}^m \text{ is a basis}.
\]

\textbf{Corollary 3.2.} To obtain ``four dimensions,'' one must first
have \(m-1=4\) or \(m=4\); neither is a consequence of the substrate
postulates.

\textbf{Lemma 3.3 (Transport rank and integer reachability).}

\[
\operatorname{im}B=H_Q, \qquad \operatorname{rank}B=m-1,
\]

\[
x\in H_Q\cap\mathbb Z^C \implies \exists f\in\mathbb Z^E,\; Bf=x.
\]

\emph{Proof.} \(B\) is totally unimodular {[}5{]}; integer reachability
follows from standard integer linear programming results.
\(\blacksquare\)

\subsection{Discrete Dirichlet
Energy}\label{discrete-dirichlet-energy}

\textbf{Lemma 3.4.} The discrete Dirichlet energy with edge weights
\(K_e\)

\[
E(f)=\sum_e K_e(f_i-f_j)^2 \ge 0, \qquad E(f)=0 \iff f \text{ is constant},
\]

i.e.~\(E=f^{\mathsf T}Lf\), \(\ker L=\operatorname{span}\{\mathbf 1\}\),
\(L(H_Q)\subseteq H_Q\), \(L|_{H_Q}\succ0\).

\textbf{Lemma 3.5 (No preference \(\Rightarrow\) uniform conductance).}
If \(\Gamma\) is edge-transitive, then \(K_e\equiv\kappa\).

\emph{Remark.} This is conditional: path graphs \(P_4\), \(P_5\) have 2
edge orbits, and no-preference does not yield uniform conductance.

\subsection{Local Quadratic Form and Spatial
Metric}\label{local-quadratic-form-and-spatial-metric}

\textbf{Lemma 3.6 (Metric-volume consistency inversion).} Continuum
approximation gives \(E\to\int Q^{ij}\partial_if\,\partial_jf\,\omega\).
Requiring metric-volume consistency

\[
\mu=\operatorname{vol}_h=\sqrt{\det h}\;\omega,
\tag{3.1}
\]

one obtains \(Q=\sqrt{\det h}\,h^{-1}\), \(\det Q=(\det h)^{d/2-1}\)
(\(d\) = spatial dimension). For \(d\neq2\) the unique inverse is

\[
h=(\det Q)^{\frac{1}{d-2}}\,Q^{-1}.
\tag{3.2}
\]

For \(d=2\), \(\det Q=1\), and \(h\mapsto\Omega^2h\) leaves \(Q\)
invariant; only the conformal class remains.

\emph{Remark.} (3.1) is the sole additional input in this step; it binds
the Dirichlet quadratic form one-to-one to the spatial metric and
removes conformal ambiguity for \(d\neq2\). The unique inverse holds for
all \(d\neq2\) and does not select any dimension.

\subsection{Splitting, Time, and
Signature}\label{splitting-time-and-signature}

\textbf{Lemma 3.7 (Reversible/irreversible splitting).}

\[
T\cong H_Q\oplus\mathbb R_\tau\cong\mathbb R^m.
\]

The regular bilinear form can only be \(g=d\tau^2-h\) or
\(g=-d\tau^2+h\). Taking the former gives
\(\operatorname{signature}(g)=(1,m-1)\).

\emph{Mechanism.} The substrate's compensatory moves are reversible
(each edge has an inverse move); step progression is irreversible.
Spatial directions carry the inverse structure; the step direction does
not. The signature originates here---there is no need to single out a
time direction among the spatial directions.

\emph{Boundary.} \(1\oplus(m-1)\) and \(h>0\) are derived; which side
carries the minus sign is a sign convention.

\subsection{Curvature and Two-Derivative
Uniqueness}\label{curvature-and-two-derivative-uniqueness}

\textbf{Lemma 3.8.} Given a smooth manifold \(M\) and nondegenerate
\(g\), the connection \(\Gamma^a{}_{bc}\), curvature \(R^a{}_{bcd}\),
Ricci \(R_{ab}\), and Einstein tensor \(G_{ab}\) are uniquely determined
by standard differential geometry {[}6{]}.

\textbf{Lemma 3.9 (Uniqueness within the two-derivative family).} Let
\(E_{ab}=\alpha R_{ab}+\beta Rg_{ab}+\gamma g_{ab}\). Using
\(\nabla^ag_{ab}=0\) and the contracted Bianchi identity
\(\nabla^aR_{ab}=\tfrac12\nabla_bR\):

\[
\nabla^aE_{ab}=\Bigl(\frac{\alpha}{2}+\beta\Bigr)\nabla_bR.
\tag{3.3}
\]

Thus when \(\nabla R\not\equiv0\), divergence-free
\(\iff\beta=-\alpha/2\),
i.e.~\(E_{ab}\in\operatorname{span}\{g_{ab},G_{ab}\}\).

\subsection{Lovelock Uniqueness and the GR
Criterion}\label{lovelock-uniqueness-and-the-gr-criterion}

\textbf{Lemma 3.10 (Lovelock uniqueness) {[}3{]}.} In 4 dimensions, a
symmetric second-rank tensor constructed from \(g\) and its derivatives
up to second order, linear in the highest derivative, and satisfying
\(\nabla^aE_{ab}=0\) must be

\[
E_{ab}=c_1G_{ab}+c_2g_{ab}.
\]

(The Gauss--Bonnet term is a topological term in 4 dimensions and does
not enter the field equations.)

\textbf{Corollary 3.11 (GR criterion).}

\[
(L)\wedge(O)\wedge(C') \implies G_{ab}+\Lambda g_{ab}=8\pi G\,T_{ab}.
\tag{3.4}
\]

Here \((C')\) requires that \(T_{ab}\) be a functional of matter data
(the metric only enters index raising/lowering and normalization) and
must not contain independent metric degrees of freedom.

\emph{Remark.} The Lovelock theorem's literal premise ``in 4
dimensions'' means it cannot be used to derive \(D=4\); dimension must
be argued separately (§6).

\subsection{Locality Bifurcation and Excluded
Routes}\label{locality-bifurcation-and-excluded-routes}

\textbf{Lemma 3.12.} The discrete Dirichlet energy is local; the
effective resistance distance
\(R_{ij}=(e_i-e_j)^{\mathsf T}L^+(e_i-e_j)\) is nonlocal.

\textbf{Lemma 3.13 (Limitations of the algebraic assembly route).} The
curvature of an equilateral complex is locked by combinatorial
structure: the deficit angles around edge \(k\) are
\(+148.41°,+77.88°,+7.36°,-63.17°\) (\(k=3,4,5,6\)), with only
\(k=3,4,5\) giving positive deficit. Algebraic assembly can therefore
only land on an equilateral, rigid subset and cannot approximate a
general metric.

{\def\LTcaptype{none} % do not increment counter
\begin{longtable}[]{@{}
  >{\raggedright\arraybackslash}p{(\linewidth - 2\tabcolsep) * \real{0.5000}}
  >{\raggedright\arraybackslash}p{(\linewidth - 2\tabcolsep) * \real{0.5000}}@{}}
\toprule\noalign{}
\begin{minipage}[b]{\linewidth}\raggedright
Excluded route
\end{minipage} & \begin{minipage}[b]{\linewidth}\raggedright
Basis
\end{minipage} \\
\midrule\noalign{}
\endhead
\bottomrule\noalign{}
\endlastfoot
Algebraic assembly (\(K\to h\)) & Unique fixed point = uniform
\(\Rightarrow\) equilateralization \(\Rightarrow\) one-parameter
rigidity \\
Effective resistance as geometry & Nonlocal; external edges break local
isotropy \\
Homogeneity degree for dimension selection & Homogeneity degree is 1 for
every \(d\); does not select dimension \\
\end{longtable}
}

\begin{center}\rule{0.5\linewidth}{0.5pt}\end{center}

\section{Metric Construction and the Lovelock
Lever}\label{metric-construction-and-the-lovelock-lever}

\subsection{Common Weight Formula}\label{common-weight-formula}

All three metric routes use the same weight formula

\[
w_{ij}=\phi_i\phi_jA_{ij},
\qquad h=\text{graph Laplacian with weight }w\text{ restricted to }H_Q,
\tag{4.1}
\]

differing only in the source of \(\phi\).

\textbf{Route 1 (closed-walk counting, \(k\to\infty\)).} The total
traversal count of edge \((i,j)\) by length-\(k\) closed walks is

\[
N^{(k)}_{ij}=k\,A_{ij}(A^{k-1})_{ij}.
\tag{4.2}
\]

After normalization, \(w_{ij}\propto\phi_i\phi_jA_{ij}\), with \(\phi\)
the Perron vector. Zero free parameters. But \(w\) depends on the Perron
eigenvector---a global quantity of the entire graph---so \((L')\) (\(g\)
is locally determined by data) fails.

\textbf{Route 2 (finite truncation \(k=n\)).} Truncating (4.2) at scale
\(k\):

\[
w^{(k)}_{ij}=\sum_{m=2}^{k} m\,A_{ij}(A^{m-1})_{ij}.
\tag{4.3}
\]

Locality holds exactly when \(k\le\) graph diameter (the effect of
distal perturbations on central edge weight ratios is exactly zero).
\(k>n\) counts unrealizable closed words; \(k<n\) misses realizable
ones, so \(k=n\) is forced.

\textbf{Route 3 (locally given \(\phi\)).} If \(\phi\) is given locally,
\((L')\) holds exactly, and all three premises are satisfied.

\textbf{Remark (Scope of the \((O)\) and \((C)\) conditions).} The
second-order \((O)\) and conservation \((C)\) checks in Routes 1--3
strictly verify that the graph Laplacian \(L\) is a second-order,
charge-conserving operator on scalar functions
(\(\ker L = \operatorname{span}\{\mathbf{1}\}\)), which guarantees that
the emergent spatial metric \(h\) is well-defined. They do not
automatically imply that the field tensor \(E_{ab}[g]\) is
divergence-free under full four-dimensional spacetime diffeomorphisms;
four-dimensional covariance \(\nabla^a E_{ab}=0\) remains an essential
macroscopic symmetry postulate registered in §7.1 (Item 8).

\subsection{Degrees of Freedom and
Parameterization}\label{degrees-of-freedom-and-parameterization}

The number of spatial metric components is

\[
\dim\operatorname{Sym}^2H_Q=\binom m2=\frac{(m-1)m}{2}.
\tag{4.4}
\]

\emph{Key distinction.} (4.4) counts spatial metric \(h\) components
(\(m=4\), 3D space: 6). Our construction \(g=d\tau^2-h\) (ADM
\(N=1,\beta=0\)) fixes \(g_{00}=1\), \(g_{0i}=0\); graph edge weights
provide only the 6 components of \(h\), never the 10 components of the
4D spacetime metric \(g_{ab}\).

\emph{Parameterization is not bijective.} The weight formula (4.1) uses
a vertex field \(\phi=(\phi_1,\dots,\phi_m)\) to give edge weights
\(w_{ij}=\phi_i\phi_jA_{ij}\). \(\phi\) has only \(m\) free parameters,
while \(\dim\operatorname{Sym}^2H_Q=\binom m2\); on \(K_m\), the
\(\binom m2\) edge weights from \(m\) vertex values always satisfy the
cross-product constraint
\(w_{12}w_{34}=w_{13}w_{24}=w_{14}w_{23}=\phi_1\phi_2\phi_3\phi_4\). It
therefore spans only a \(\le m\)-dimensional submanifold of
\(\operatorname{Sym}^2H_Q\), not a bijection. The injectivity of
\(w\mapsto h\) (for independent weights) holds, but the \(\phi\)
parameterization of (4.1) is a local parameterization on a submanifold,
not proven to cover the full space.

\emph{Graviton degrees of freedom.} In general 4D GR: 10 metric
components \(-4\) coordinate gauge \(-4\) initial-value constraints
\(=2\) physical propagating degrees of freedom,
i.e.~\(\frac{D(D-3)}2=2\).

\subsection{Refinement Convergence and Nonlocality of the
Limit}\label{refinement-convergence-and-nonlocality-of-the-limit}

The edge-weight ratio converges with truncation \(k\), but extremely
slowly, and the limit is the Perron form \(P_1=vv^{\mathsf T}\),
globally nonlocal. The deviation decay length is

\[
\tau\approx\frac{2N^2}{3\pi^2}.
\tag{4.5}
\]

\emph{Conclusion.} The effective metric has a limit under refinement;
the limit is the Perron form, nonlocal \(\Rightarrow\) there exists no
UV-fixed local metric.

\begin{center}\rule{0.5\linewidth}{0.5pt}\end{center}

\section{Continuum Limit and Causal
Structure}\label{continuum-limit-and-causal-structure}

\subsection{\texorpdfstring{5.1
\(\Gamma\)-Convergence}{5.1 \textbackslash Gamma-Convergence}}\label{gamma-convergence}

Lovelock requires \(\Gamma\)-convergence (local structure convergence),
not a UV fixed point (self-similarity) {[}7{]}. A standard fact of
homogenization theory: non-self-similar coefficient families still
\(\Gamma\)-converge. With fixed step count \(k=8\), the profile maximum
difference between adjacent \(N\) converges at approximately \(O(a^2)\).

\emph{Boundary.} \(\Gamma\)-convergence is not proven (no compactness,
no \(\Gamma\)-liminf); only convergence and order are measured, and only
on a 1D chain with a single profile.

\subsection{Causal Structure: Support Cone and Cone
Edge}\label{causal-structure-support-cone-and-cone-edge}

Finite causal velocity requires saturation (per-point capacity), not
current relaxation. Out-of-cone signals are exponentially small but
nonzero pulled fronts {[}8{]}.

\textbf{The support cone is always exact}: nearest-neighbor transport
expands at most one lattice site per step,
\(\operatorname{supp}\rho_n\subseteq[-n,n]\), independent of \(B\).

\textbf{Cone-edge visibility} is derived analytically: on the
\(e^{-\mu x}\) mode, the one-step mapping characteristic equation is

\[
\Lambda^2=B\cosh^2\mu \implies \lambda(\mu)=\tfrac12\log B+\log\cosh\mu.
\tag{5.1}
\]

The cone-edge density decay rate is

\[
\varepsilon(B)=\tfrac12\log\Bigl(\tfrac4B\Bigr).
\tag{5.2}
\]

\emph{Visible cone edge \(\iff B=4\).} \(B=4\) is upgraded from a model
parameter to a self-consistency condition.

\begin{center}\rule{0.5\linewidth}{0.5pt}\end{center}

\section{Dimension Account}\label{dimension-account}

\textbf{This is the weakest link of the entire chain.}

\subsection{The Core Postulates Do Not Constrain Dimension
(No-Go)}\label{the-core-postulates-do-not-constrain-dimension-no-go}

\textbf{Proposition 6.1 (Classification-level no-go).} Let
\(\mathcal A\) be the substrate core postulates. Then for every integer
\(m\ge2\), there exists a model \(\mathcal M_m\) of \(\mathcal A\) with
channel-set cardinality \(|C|=m\).

\emph{Proof.} Take \(C=\{0,\dots,m-1\}\),
\(E=\{\{i,i+1\}:i\in\mathbb Z_m\}\) (cycle graph \(C_m\)). Verify
postulate by postulate: \(C_m\) is finite and connected; every move has
an inverse; all \(2m\) directed moves are taken with equal weight; step
index progresses irreversibly; closed state \(x=0\) exists. Numerical
verification (\(m=2,\dots,12\)): \(\operatorname{rank}B=m-1\),
\(\ker L=\operatorname{span}\{\mathbf1\}\), \(L|_{H_Q}\succ0\). \(C_m\)
is edge-transitive, so the no-preference \(\Rightarrow\) uniform
conductance premise also holds. \(\blacksquare\)

\[
\implies \text{any statement depending only on }D\text{ cannot be derived from the substrate core postulates}.
\tag{6.1}
\]

\subsection{Excluding Low
Dimensions}\label{excluding-low-dimensions}

{\def\LTcaptype{none} % do not increment counter
\begin{longtable}[]{@{}
  >{\centering\arraybackslash}p{(\linewidth - 4\tabcolsep) * \real{0.4000}}
  >{\raggedright\arraybackslash}p{(\linewidth - 4\tabcolsep) * \real{0.3000}}
  >{\raggedright\arraybackslash}p{(\linewidth - 4\tabcolsep) * \real{0.3000}}@{}}
\toprule\noalign{}
\begin{minipage}[b]{\linewidth}\centering
\(D\)
\end{minipage} & \begin{minipage}[b]{\linewidth}\raggedright
Conclusion
\end{minipage} & \begin{minipage}[b]{\linewidth}\raggedright
Basis
\end{minipage} \\
\midrule\noalign{}
\endhead
\bottomrule\noalign{}
\endlastfoot
\(2\) & No gravitational dynamics &
\(R_{ab}=\tfrac12Rg_{ab}\implies G_{ab}\equiv0\) \\
\(3\) & No graviton & Weyl tensor vanishes identically; curvature fully
determined by local matter \\
\(4\) & Propagating graviton exists & see §6.3 \\
\end{longtable}
}

\subsection{\texorpdfstring{6.3 Uniqueness of \(D=4\) from Graviton
Degrees of
Freedom}{6.3 Uniqueness of D=4 from Graviton Degrees of Freedom}}\label{uniqueness-of-d4-from-graviton-degrees-of-freedom}

The \(D\)-dimensional graviton degrees of freedom are
\(\frac{D(D-3)}2\):

{\def\LTcaptype{none} % do not increment counter
\begin{longtable}[]{@{}lrrrrrr@{}}
\toprule\noalign{}
\(D\) & 2 & 3 & 4 & 5 & 6 & 7 \\
\midrule\noalign{}
\endhead
\bottomrule\noalign{}
\endlastfoot
\(\frac{D(D-3)}2\) & \(-1\) & \(0\) & \(\mathbf2\) & \(5\) & \(9\) &
\(14\) \\
\end{longtable}
}

\[
\frac{D(D-3)}2=2 \iff D^2-3D-4=0 \iff D\in\{-1,4\} \implies D=4.
\tag{6.2}
\]

\emph{Screening chain.}
\(\{2,3,\dots\}\xrightarrow{\ D(D-3)/2>0\ }\{4,5,\dots\}\xrightarrow{\text{minimality}}4\).

\emph{Boundary.} Minimality, even-dimension reading, and
sign--polarization identification are not consequences of the substrate
postulates.

\subsection{Candidate Constraints as Synonymous
Reformulations}\label{candidate-constraints-as-synonymous-reformulations}

Of 15 candidate constraints tested item by item, items 1--10 all have
survival set \(\{4\}\); they are all synonymous reformulations of
\((Z_2)'\) and do not constitute independent evidence.

Among them, item 7, \(b_1(K_m)=m-1\), is examined in detail below.

\subsection{Representation-Theoretic Status of the
Cycle--Displacement
Duality}\label{representation-theoretic-status-of-the-cycledisplacement-duality}

As an \textbf{oriented} cycle space, \(Z_1(K_4)\) decomposes as the
standard \(S_4\)-module \((2,1,1)\cong(3,1)\otimes\mathrm{sgn}\), which
is the representation of axial vectors (pseudovectors, flux, magnetic
field). The displacement space is \(H_Q=(3,1)\), the standard
polar-vector representation. The two are distinct \(S_4\)-modules (they
differ by the sign representation). \textbf{However, in
three-dimensional geometric realization} (\(K_4\) embedded as a regular
tetrahedron in \(\mathbb R^3\)), the Hodge star operator

\[
\star:\Lambda^2\mathbb R^3\longrightarrow\mathbb R^3
\]

establishes a canonical duality between the axial cycle space and the
polar displacement space. This is not a numerical coincidence: it is the
group-theoretic counterpart of §6.8(ii) ``antisymmetric composition
closure''---the identity
\(\Lambda^2\mathbb R^3\cong\mathbb R^3\otimes\mathrm{sgn}\) is precisely
the Hodge duality between closed 2-forms (loops, flux) and 1-forms
(displacements) in three dimensions {[}9{]}.

\subsection{Dimension-Blind No-Go}\label{dimension-blind-no-go}

All four metric-side structures \(u=h/\sqrt n\), \(\mathrm H(w)\),
\(\min(k,n-k)\), \(T(M,n)\) do not reference the adjacency matrix \(A\)
and take the same value on any graph.

\[
\implies \textbf{the metric side cannot exclude low dimensions (e.g., }d=1\text{)}.
\tag{6.3}
\]

\subsection{The Closure--Motion Correspondence
Principle}\label{the-closuremotion-correspondence-principle}

\textbf{The core postulates do not constrain dimension} (Proposition
6.1). What fixes the dimension is this correspondence:

\begin{quote}
\textbf{Closure--motion correspondence {[}named input{]}.} The axiom,
read at the counting level: \(\dim\ker B=\dim\operatorname{im}B\)
(degrees of freedom of closed loops = degrees of freedom of motion).
\end{quote}

Here \(B\) is the incidence matrix of compensatory moves:
\(\operatorname{im}B=H_Q\) is the zero-sum state, and \(\ker B\) is the
closed-loop space.

\begin{align}
&\text{dictionary}: \quad d=m-1, \quad \text{uniquely fixed by } \binom m2-(m-1)=\binom{m-1}2;\\
&\text{reading}: \quad \binom m2-(m-1)=m-1 \iff m=4 \quad(\text{unique});\\
&\text{no preference (edge-transitive)}: \quad \deg\text{ equal} \implies t\,m=4(m-1) \iff (\deg,m)=(3,4).
\end{align}

\textbf{Machine check}: on \(K_m\), \(\dim\ker B=\binom m2-(m-1)\) and
\(\dim\operatorname{im}B=m-1\) are equal only when \(m=4\) (\(m=2\)
gives \(0\ne1\), \(m=3\) gives \(1\ne2\), \(m\ge5\) always gives
\(\dim\ker B>\dim\operatorname{im}B\)).

Together the three give: \textbf{spatial dimension \(d=3\), channel
count \(m=4\), \(\deg=3\), spacetime dimension \(D_{\rm st}=d+1=4\)}.

\[
D_{\rm st}=4\text{ follows from two independent routes: the substrate-side premise is the closure--motion correspondence, the physical-side premise is "propagating gravitons exist".}
\tag{6.5}
\]

\emph{Why the correspondence cannot be a consequence of the postulates.}
The substrate gives \(\Gamma=(C,E)\) together with ``motion''; the
postulates only constrain \emph{how} to move---``full branching'' means
every \emph{allowed} motion is instantiated, and does not constrain
which pairs have allowed moves. Any condition of the form ``\(\Gamma\)
must be complete/self-dual'' is a choice of \(\Gamma\), not a
consequence of the postulates. Counterexample: cycle graph \(C_m\)
(\(m\ge4\)) satisfies the postulates and edge-transitivity, but
\(\dim\ker B=1\ne m-1\).

\subsection{Equivalent Formulations of the Correspondence
Principle}\label{equivalent-formulations-of-the-correspondence-principle}

The counting content \(\binom d2=d\) has three equivalent formulations:

\textbf{(i) Pairing uniqueness.} Each edge has exactly one opposite edge
\(\iff\binom{m-2}2=1\iff m=4\).

\textbf{(ii) Antisymmetric composition closure.}
\(\dim\Lambda^2V=\dim V\iff D=3\): the ``relative relation'' of two
directions is again a direction (cross-product structure). I.e.
\(\binom D2=D\iff D\in\{0,3\}\), with \(D=3\) the unique nontrivial
solution.

\textbf{(iii) \(A_3\) self-duality.} The 6 edge directions
\(\pm(e_i-e_j)\) of \(K_4\) are exactly the 6 positive roots of \(A_3\);
\(\operatorname{Aut}(K_4)=S_4\) is the Weyl group of \(A_3\). Among
\(A_n\), only \(A_3\) satisfies ``number of relations = dimension'':
\(\binom{n+1}2-n=n\iff n=3\). This is precisely the accidental
isomorphism \(A_3\cong D_3\) (\(\mathfrak{so}(3)\cong\mathbb R^3\)
cross-product isomorphism).

\[
\implies \Gamma=K_4:\ 4\text{ vertices},\ \deg=3,\ 6\text{ edges},\ 3\text{ opposite-edge pairs};\ \text{angle}=\arccos(-1/3)=109.4712°;\ \operatorname{Aut}(K_4)=S_4.
\tag{6.6}
\]

\begin{center}\rule{0.5\linewidth}{0.5pt}\end{center}

\section{Axiomatic Budget}\label{axiomatic-budget}

\subsection{What the Lovelock Lever
Presupposes}\label{what-the-lovelock-lever-presupposes}

{\def\LTcaptype{none} % do not increment counter
\begin{longtable}[]{@{}
  >{\raggedleft\arraybackslash}p{(\linewidth - 4\tabcolsep) * \real{0.3333}}
  >{\raggedright\arraybackslash}p{(\linewidth - 4\tabcolsep) * \real{0.3333}}
  >{\raggedright\arraybackslash}p{(\linewidth - 4\tabcolsep) * \real{0.3333}}@{}}
\toprule\noalign{}
\begin{minipage}[b]{\linewidth}\raggedleft
\#
\end{minipage} & \begin{minipage}[b]{\linewidth}\raggedright
Presupposition
\end{minipage} & \begin{minipage}[b]{\linewidth}\raggedright
Nature
\end{minipage} \\
\midrule\noalign{}
\endhead
\bottomrule\noalign{}
\endlastfoot
1 & \(D=4\) (literal hard-coding) & Cannot select dimension via this
lever \\
2 & Metric exists and is smooth & Input \\
3 & Metric-volume consistency \(\mu=\operatorname{vol}_h\) & Sole
additional input of (3.1) \\
4 & Second-order \((O)\) & Excludes higher-order terms \\
5 & Conserved source \((C')\) & \((C)\) as-is is vacuous \\
6 & Source = relativistic matter & Additional claim \\
7 & Metric nondegenerate and Lorentzian & Input \\
8 & 4D diffeomorphism covariance \(\nabla^aE_{ab}=0\) & Input: the
preferred foliation of \(g=d\tau^2-h\) gives only 3D spatial
diffeomorphism invariance \\
\end{longtable}
}

\subsection{Absence of Intrinsic
Scales}\label{absence-of-intrinsic-scales}

\[
\text{The dimensionful degrees of freedom in the geometric sector are exactly }1.
\tag{7.1}
\]

\emph{Evidence.} The layer coefficient (edge \((0,1)\))
\(\binom{m-1}{m/2}\) is an integer (\(m=2,4,6,8\) gives \(1,3,10,35\)).
All raw metric-side data are integers \(\Rightarrow\) they cannot carry
length.

\emph{Unit-transformation invariance.} Under \(W\to sW\) (\(s\) spanning
8 orders of magnitude), normalized eigenvalue deviation is
\(1.89\times10^{-15}\), edge-weight ratio deviation is \(0\).

\emph{On \(G,\Lambda\).} Dimensional constants are units, not parameters
{[}10{]}. ``Deriving the numerical values of \(G,\Lambda,\hbar\)'' is
therefore not a well-posed problem.

\begin{center}\rule{0.5\linewidth}{0.5pt}\end{center}

\section{Main Results and
Boundaries}\label{main-results-and-boundaries}

\subsection{Main Result}\label{main-result}

\[
\begin{aligned}
&\text{substrate (zero-sum constraint) + (inputs: embedding / source class / scale)}\\
&\quad\Longrightarrow (L)\wedge(O)\wedge(C')\\
&\quad\Longrightarrow G_{ab}+\Lambda g_{ab}=8\pi G\,T_{ab}\quad(\text{in }D=4).
\end{aligned}
\tag{8.1}
\]

\[
\text{the shape of the field equations is uniquely forced; the metric and dimension are inputs; the source is a no-go.}
\tag{8.2}
\]

\subsection{Layering of ``Dimension Is an
Input''}\label{layering-of-dimension-is-an-input}

``Input'' here refers specifically to the Lovelock layer: the shape
theorem \(E_{ab}=\alpha G_{ab}+\beta g_{ab}\) takes \(D\) as an
argument, and Lovelock itself does not select dimension. But the value
\(D=4\) is not a bare free input; it is given by the readout chain: peak
condition locks \(L=4\) (Gate 1) + closure-equality
\(\dim\ker B=\dim\operatorname{im}B\) locks \(K_4/A_3\). This readout
chain is half-locked: the peak condition hard-presses to \(L=4\); the
step ``\(L=4\) is exactly the physical spacetime dimension'' is only
count-level
equality---\(Z_1(K_4)\cong(2,1,1)\cong(3,1)\otimes\mathrm{sgn}\) and
\(H_Q=(3,1)\) are distinct \(S_4\)-modules differing by a parity sign,
but in three-dimensional geometry form an axial--polar vector duality
via the Hodge star.

\subsection{Boundaries That Must Be
Stated}\label{boundaries-that-must-be-stated}

\begin{enumerate}
\def\labelenumi{\arabic{enumi}.}
\tightlist
\item
  \textbf{Division of labor of Lovelock}: it only answers ``what must
  the field equations look like,'' not ``why does geometry have
  dynamics.''
\item
  \textbf{Local \(K_4\) fiber vs.~macroscopic grid are two structures.}
  \(\dim H_Q=m-1=3\) holds only for the \emph{single-point} local
  \(K_4\) fiber; on the macroscopic \(N\)-vertex grid,
  \(\dim H_Q=N-1\ne3\). The construction bridge ``how local \(K_4\)
  fibers/cells tile a macroscopic 3D manifold grid'' is missing; it is
  an open item.
\item
  \textbf{\(L=4\) uses Gate 1 only.} The peak condition locks \(L=4\);
  this \(k=2\) partition does not pass the contextuality gate or the
  density gate, so it serves only as dimension readout, not as a
  quantumness criterion.
\end{enumerate}

\begin{center}\rule{0.5\linewidth}{0.5pt}\end{center}

\section{Open Problem}\label{open-problem}

\textbf{Continuum limit (\(\Gamma\)-convergence)} is unproven: no
compactness, no \(\Gamma\)-liminf; only on a 1D chain with a single
profile. This is analytic (compactness / variational limit) and requires
independent work in homogenization theory.

\begin{center}\rule{0.5\linewidth}{0.5pt}\end{center}

\section{Conclusion}\label{conclusion}

\[
\begin{aligned}
&\text{substrate (zero-sum constraint)}\\
&\quad\Longrightarrow \dim H_Q=m-1;\ \operatorname{im}B=H_Q;\ \text{discrete Dirichlet energy};\\
&\quad\Longrightarrow\text{(continuum + metric-volume consistency) spatial metric }h;\\
&\quad\Longrightarrow\text{(reversible/irreversible splitting) signature }(1,m-1);\\
&\quad\Longrightarrow\text{(locality}\wedge\text{second-order}\wedge\text{conserved source)}\\
&\quad\Longrightarrow G_{ab}+\Lambda g_{ab}=8\pi G\,T_{ab}\quad(\text{in }D=4).
\end{aligned}
\tag{10.1}
\]

\textbf{Shape is unconditional; entry is conditional}:
\((L)\wedge(O)\wedge(C')\Rightarrow\) Einstein form is unconditional
(Lovelock); dimension, metric, and source are three entries, all
non-consequences of the postulates.

Current progress is \textbf{conditional recovery}, not
\textbf{unconditional derivation}. To reach \emph{derivation}, the same
ledgers are missing: the premise of dimension, the origin of the metric,
and the price tag of the source.

\begin{center}\rule{0.5\linewidth}{0.5pt}\end{center}

\section*{References}\label{references}

{[}1{]} R. Penrose, \emph{The Road to Reality}, Vintage Books, 2005.

{[}2{]} C. J. Isham, ``Structural issues in quantum gravity,''
arXiv:gr-qc/9510034, 1995.

{[}3{]} D. Lovelock, ``The Einstein tensor and its derivatives,''
\emph{J. Math. Phys.} \textbf{12}, 498 (1971).

{[}4{]} W. Feller, \emph{An Introduction to Probability Theory and Its
Applications}, Vol. 1, 3rd ed., Wiley, 1968.

{[}5{]} A. Schrijver, \emph{Theory of Linear and Integer Programming},
Wiley, 1986.

{[}6{]} M. P. do Carmo, \emph{Riemannian Geometry}, Birkhäuser, 1992.

{[}7{]} D. G. Cioranescu, P. Donato, \emph{An Introduction to
Homogenization}, Oxford University Press, 1999.

{[}8{]} J. D. Murray, \emph{Mathematical Biology}, Vol. II: Spatial
Models and Biomedical Applications, 3rd ed., Springer, 2003.

{[}9{]} H. Weyl, \emph{Classical Groups: Their Invariants and
Representations}, Princeton University Press, 1939.

{[}10{]} M. J. Duff, L. B. Okun, G. Veneziano, ``Trialogue on the number
of fundamental constants,'' \emph{JHEP} \textbf{0203}, 023 (2002),
arXiv:physics/0110060.

\end{document}
